\section{Conclusion}
\label{sec:conclusion}

We added a single margin, an endogenous brown/green durable adoption
choice, to an otherwise standard small open economy HANK model, and asked how
crisis policy interacts with it. The central result is a non-neutrality: a retail
price cap stabilises output but, by holding the brown price fixed, switches the
green adoption margin off, whereas a transfer of the same fiscal cost delivers
comparable stabilisation and leaves the margin intact. Shielding consumers from
the relative price is not neutral for the energy transition, and Europe's own
heat-pump cycle (boom when gas was expensive, bust when subsidised gas made it
cheap again) shows the margin operating in the data. The same logic, applied
ex ante, yields the second result: a permanent carbon tax that greens the steady
state provides no crisis insurance to first order, and the insurance value does
not grow as the shock becomes persistent, because pre-greening spends an adoption
absorber that works best when the economy enters the crisis with room to green.
The case for carbon pricing is the externality; the case for crisis protection is
income support, paid flat.

\label{sec:limits}%
The model is built to isolate one margin, and its limitations indicate the next
steps. The taste scale $\sigma_\varepsilon$ is calibrated to the
quasi-experimental vehicle-adoption literature rather than estimated on
European durables; the identification in
Section~\ref{sec:taste_id} makes the mapping explicit, and estimating the
adoption elasticity directly is the next step; heat-pump sales against
the electricity--gas price ratio across European countries would be the natural
design. The switching cost is booked as an import by accounting
necessity under the inherited markup structure; the ordering of the policy
results does not depend on this convention, but the sign of the adoption
channel's output contribution does, and we have noted it as conditional
throughout (Appendix~\ref{app:booking}). The green side of the economy is an accounting construct
rather than a technology: there is no green supply sector with capacity and
investment, and no installer income distribution beyond a lump-sum rebate. The
supply side more broadly has no capital, which likely makes the transfer
somewhat more expansionary and the supply shock more inflationary than
they would be with an investment margin. Finally, energy demand is homothetic
conditional on the durable type, so the model is silent on the
energy-share-of-the-poor mechanism of \citet{bayer2026hank}; the two
distributional channels complement each other.

None of these caveats changes what the paper establishes: the extensive margin
of green adoption is a quantitatively important channel of energy-crisis policy,
strong enough that the choice between shielding prices and shielding incomes
is also a choice about the speed of the transition.
